% !TEX root = ../Projektdokumentation.tex
\section{Implementierungsphase}
\label{sec:Implementierungsphase}

Die Implementierung umfasst die physische Vorbereitung des Servers,
die Installation und Konfiguration des \ac{PBS} sowie die Anbindung
an alle beteiligten Systeme. Abbildung~\ref{fig:Architektur} zeigt
die Gesamtarchitektur der Referenzkonfiguration im Überblick.

\begin{figure}[htbp]
    \centering
    \includegraphics[width=0.95\linewidth,keepaspectratio]{architecture}
    \caption{Gesamtarchitektur der PBS-Referenzkonfiguration:
             PVE-Cluster, PBS odcf-fs04, LDAP, Prometheus und
             bestehende TSM-/Tape-Sicherung}
    \label{fig:Architektur}
\end{figure}

\subsection{Beschaffung und Vorbereitung}
\label{sec:Beschaffung}
Da als Hardware ein vorhandener IBM System x3755 M3 nachgenutzt wird,
entfällt eine externe Beschaffung. Der Server wird in DCODCF-Rack14
bereitgestellt und mit einem Hardware-\ac{RAID}-Controller
konfiguriert: \ac{RAID}\,1 für das Betriebssystem und
\ac{RAID}\,5 über 6$\times$4\,TB HDDs (Modell ST4000NM033) für den
Datastore (ca.~20\,TB nutzbar).
Zusätzlich wird eine \ac{SFP}-Netzwerkkarte mit 2$\times$10\,GbE
eingebaut und per \ac{LACP} zu einem logischen Interface gebündelt;
die Verkabelung läuft über die Patchpanels in Rack14/16 zum
Cisco-Switch \texttt{aci-dc03-og2-r16-aly-2153}. Die
Netzwerk-Infrastruktur (\emph{server}-\ac{VLAN},
Firewall-Freischaltungen) ist bereits vorhanden.
Abschließend wird das aktuelle \ac{PBS}-4-ISO heruntergeladen und
als bootfähiges Installationsmedium vorbereitet.

Auf Linux-Seite wird das Bond als \texttt{bond0} im Modus
\texttt{802.3ad} (\ac{LACP}) mit \texttt{layer3+4}-Hashing
konfiguriert~\citep{IEEE802ad}; der vollständige Auszug aus
\texttt{/etc/network/interfaces} findet sich in
Listing~\ref{lst:interfaces} im Anhang.


\subsection{Installation und Grundkonfiguration}
\label{sec:InstallationPBS}
Die Installation erfolgt direkt vom \ac{PBS}-4-ISO. Der geführte
Installer richtet Debian mit dem \ac{PBS}-Paket auf dem
\ac{RAID}\,1-System-Volume ein. Anschließend wird der Host wie
folgt konfiguriert: Hostname \texttt{odcf-fs04} mit der IP
\texttt{10.131.196.78} im \emph{server}-\ac{VLAN}, Eintrag im
\ac{DKFZ}-DNS, Konfiguration des
\texttt{pbs-no-subscription}-Repositories (keine
Enterprise-Subscription) und vollständiges \texttt{apt}-Upgrade.
Die NTP-Synchronisation läuft gegen die internen Zeitserver des
\ac{DKFZ}. Die Web-Oberfläche des \ac{PBS} ist unter
\url{https://10.131.196.78:8007} erreichbar.

Tatsächlicher Zeitaufwand für Installation und Grundkonfiguration:
ca.~fünf Stunden. Der Mehraufwand gegenüber den geplanten vier
Stunden resultiert aus den Hardware-\ac{RAID}-Komplikationen beim
initialen Disk-Einrichtung.


\subsection{Konfiguration des Datastores}
\label{sec:Datastore}
Auf dem \ac{RAID}\,5-Volume wird ein \texttt{ext4}-Dateisystem
angelegt und unter \texttt{/mnt/datastore/backup} eingehängt.
Darauf legt der \ac{PBS} den Datastore \texttt{backup}
an~\citep{PBSAdminGuide}; die Konfiguration ist in
Listing~\ref{lst:datastore} im Anhang dokumentiert. Der Ablauf eines
Backup-Jobs vom Scheduler-Trigger über die Chunk-Übertragung bis zum
Verify ist in Abbildung~\ref{fig:BackupSeq} dargestellt.

\begin{figure}[htbp]
    \centering
    \includegraphics[width=0.85\linewidth,keepaspectratio]{backup-sequence}
    \caption{Ablauf eines Backup-Jobs: PVE-Hypervisor authentifiziert
             sich am PBS, überträgt Chunks inkrementell und löst
             anschließend den Prune aus}
    \label{fig:BackupSeq}
\end{figure}

Auf dem Datastore wird zusätzlich ein wöchentlicher Verify-Job
eingerichtet, der die Prüfsummen aller gespeicherten Chunks gegen
die Metadaten vergleicht und so die Integrität langfristig
sicherstellt.


\subsection{Anbindung an das DKFZ-LDAP}
\label{sec:LDAPAnbindung}
Für die Anmeldung regulärer Administratoren wird der \ac{PBS} an
das \ac{DKFZ}-\ac{LDAP} per LDAPS auf Port~636 angebunden.
Berechtigt sind ausschließlich Benutzer der \ac{AD}-Gruppe
\texttt{ODCF-SysAdmins}; nach dem Login an der Web-Oberfläche
können sie Backup-Jobs überwachen, Restores starten und den
Datastore verwalten. Der lokale \texttt{root}-Account bleibt für
die Initialkonfiguration und Notfallzugriffe erhalten.
Die Realm-Konfiguration ist in Listing~\ref{lst:ldap} im Anhang
abgebildet; vertrauliche DKFZ-spezifische Bind-Daten sind dort
geschwärzt.

Tatsächlicher Zeitaufwand: ca.~30~Minuten.


\subsection{Integration in den PVE-Cluster}
\label{sec:PVEIntegration}
Für die Anbindung des \ac{PBS} an den \ac{PVE}-Cluster wird der
lokale \texttt{root}-Account des \ac{PBS} verwendet. Im \ac{PVE}
wird der \ac{PBS} als Storage vom Typ \emph{Proxmox Backup
Server}~\citep{PVEAdminGuide} hinzugefügt; dabei werden
Server-Adresse (\texttt{odcf-fs04}), Datastore-Name,
\texttt{root}-Zugangsdaten sowie der TLS-Fingerprint hinterlegt.
Ein Test-Backup einer unkritischen \ac{VM} verifiziert die korrekte
Anbindung, bevor produktive Jobs aufgesetzt werden.

Tatsächlicher Zeitaufwand: ca.~30~Minuten.


\subsection{Einrichtung der Backup-Jobs}
\label{sec:BackupJobs}
Die produktiven Backup-Jobs werden im \ac{PVE} pro \ac{VM}-Klasse
angelegt. Sowohl für kritische als auch für Test-\acp{VM} sind zwei
Sicherungen pro Tag (00:00 und 12:00~Uhr) mit der Prune-Regel
\texttt{keep-last=3} konfiguriert.
Jeder neu angelegte Job wird zunächst manuell ausgelöst, um
Laufzeiten und Bandbreitenbedarf zu beobachten, und erst danach
auf den automatischen Zeitplan umgestellt.


\subsection{Monitoring und Härtung}
\label{sec:Monitoring}
Für Fehlerbenachrichtigungen wird der eingebaute E-Mail-Mechanismus
des \ac{PBS} über das SMTP-Relay des \ac{DKFZ} konfiguriert.
Zusätzlich wird auf \texttt{odcf-fs04} der
\texttt{natrontech/pbs-exporter}~v0.8.0~\citep{PBSExporter} als
systemd-Unit installiert, der über Port~9100 Metriken
(Datastore-Belegung, Backup-Alter je \ac{VM}, Verify-Status) an
das zentrale Prometheus der \ac{ODCF} liefert. Die Authentifizierung
erfolgt über einen dedizierten \ac{API}-Token
(\texttt{root@pam!pbs-exporter}, Rolle \texttt{Audit}).

Zur Härtung des Hosts wird \texttt{firewalld}~\citep{FirewalldDocs} eingesetzt: Aktiv
sind ausschließlich \texttt{ssh} (22/tcp), der \ac{PBS}-Port
(8007/tcp) und der Exporter-Port (9100/tcp); alle übrigen
Default-Dienste (cockpit, http, https) sind entfernt. Der Ablauf
eines File-Restores, vom Login an der Web-Oberfläche bis zur
Datei-Extraktion aus einem Snapshot, ist in
Abbildung~\ref{fig:RestoreSeq} dargestellt.

\begin{figure}[htbp]
    \centering
    \includegraphics[width=0.85\linewidth,keepaspectratio]{restore-sequence}
    \caption{File-Restore: Admin startet über die PVE-Weboberfläche
             eine temporäre Restore-VM, browst im Snapshot und
             extrahiert die gewünschte Datei}
    \label{fig:RestoreSeq}
\end{figure}

